% !TEX root = ../main.tex

\chapter{真机实验}

\section{本章定位与目标}

如果说第四章主要回答“系统在仿真后端能否形成闭环”，那么第五章要继续回答的问题就是“这套闭环系统能否真正进入 Orin 与 Lite3 真机平台，并在真实感知、真实通讯和真实运动约束下保持有效”。因此，本章的重点不再是重新论证 NoMaD 的算法原理，而是从部署流程、平台接口、运行模式、数据记录与安全机制等维度，对真机实验进行系统梳理。

需要强调的是，真机实验在本文中并不等价于“把代码拷过去运行一次”。对足式平台而言，真正的部署工作至少包含四类内容：首先，高层 NoMaD 推理必须在边缘计算设备上稳定加载和推理；其次，运动控制桥接必须能够把高层 waypoint 转化为底层可接受的速度命令；再次，系统应提供导航模式、探索模式以及多目标任务的组织入口；最后，实验过程还应具备可视化、图像保存、日志记录和安全停机能力。也正是在这个意义上，本章讨论的是“真机系统能否成立”，而不仅仅是“真机能否动起来”。

\begin{figure}[htbp]
    \centering
    \fbox{\rule{0pt}{58mm}\rule{0.86\linewidth}{0pt}}
    \caption{图位占位：Orin 与 Lite3 真机部署总体架构图（待补）}
    \label{fig:real-overview-placeholder}
\end{figure}

\section{Orin 真机部署架构}

\subsection{边缘计算平台的角色}

在本文部署方案中，Jetson Orin 承担高层导航主机的角色，即在其上完成相机读取、NoMaD 模型加载、扩散采样推理、状态机调度、导航主机交互和图像记录等工作。之所以选择 Orin 这类边缘计算平台，是因为它能够在机载条件下提供相对完整的 GPU 推理能力，从而使高层视觉导航不必完全依赖远程工作站。

但与此同时，Orin 也带来了明确约束。与第三章中使用 RTX 4090 等设备进行离线训练和统计实验不同，Orin 的算力与显存更加有限，因此高层推理、图像显示、图像保存和控制桥接必须争夺同一套本地资源\cite{ddim2020,nomad2023}。这意味着真机部署必须从一开始就区分“调试模式”和“正式运行模式”：调试时允许开启可视化和更多日志，正式运行时则优先保障闭环控制稳定性。

\subsection{软件层次与控制链路}

当前 Orin 真机部署沿用第四章的分层思想，但进一步落实为四个可操作的软件层次：高层为 NoMaD 推理模块，中层为 waypoint 到速度参考的控制映射，低层为 Lite3 UDP 控制桥接，底层则为 Lite3 运动主机与步态控制系统。这一链路可以概括为

\begin{equation}
\text{Camera} \rightarrow \text{NoMaD Inference} \rightarrow \text{Waypoint-to-Cmd} \rightarrow \text{UDP Bridge} \rightarrow \text{Lite3}.
\end{equation}

需要进一步说明的是，真机相机并不要求原始输出分辨率与视觉编码器输入尺寸完全一致。当前部署链路会在推理前自动将图像缩放到策略配置中的 `image\_size`，例如 EfficientNet-B0、ConvNeXt-Tiny 与 ResNet-50 默认使用 $96\times96$ 输入，DINOv2-Small 使用 $98\times98$ 输入。因此，对真机部署更关键的要求不是“相机必须直接输出编码器尺寸”，而是“预处理过程与训练配置保持一致”。在此基础上，相机标定应被视为可选增强步骤而非绝对前置条件：对于普通镜头且畸变不明显的相机，可以先不标定；而对于广角镜头、鱼眼镜头或边缘直线明显弯曲的相机，则建议在桥接层加入去畸变处理，以减小真实图像分布与训练分布之间的偏移。

需要特别说明的是，Orin 并不是依赖额外桌面显卡工作的“小型 PC”，而是自带 NVIDIA 集成 GPU 的 Jetson 边缘平台。因此，本文所说的“Orin 上使用 CUDA 推理”，指的是利用其片上集成 GPU 完成 PyTorch 推理，而不是额外接入 RTX 系列独立显卡。从工程实现上看，决定 Orin 能否承担真机推理任务的关键条件，不是“是否存在桌面级 NVIDIA 显卡”，而是当前 JetPack、CUDA 运行时与 Jetson 版 PyTorch 是否正确匹配，使得集成 GPU 能够在程序中被识别为 `cuda` 设备。

这种结构的优势在于：高层 NoMaD 仍保持与仿真阶段一致的策略输入输出定义，而低层机器人通讯和步态细节则被吸收到桥接层与运动主机之中。因此，从系统工程角度看，第五章并不是另起一套真机专用算法，而是在第四章统一闭环系统的基础上切换到真实后端。

\subsection{为何需要真机桥接层}

Lite3 的运动主机并不直接接受“扩散轨迹”或“视觉目标图像”这类高层抽象概念，它接受的是速度指令、心跳信息和模式切换命令。因此，在 Orin 和机器人之间必须存在一个真机桥接层，用来完成三项关键任务：其一，把高层 motion command 映射为底层速度控制命令；其二，维护与运动主机之间的 UDP 通讯与心跳；其三，对外暴露统一平台接口，使导航主机与状态机能够在不理解底层协议细节的前提下驱动真机。

这层桥接设计与第四章中的平台抽象接口是一致的。它既是“同一高层架构、双后端复用”在工程上的落地形式，也是真机部署安全性的关键保障。没有桥接层，导航主机将不得不直接操作底层通讯协议，这会使系统既难以维护，也难以在不同硬件平台之间扩展。

\begin{figure}[htbp]
    \centering
    \fbox{\rule{0pt}{55mm}\rule{0.86\linewidth}{0pt}}
    \caption{图位占位：Orin 推理、桥接层与 Lite3 运动主机数据流图（待补）}
    \label{fig:real-bridge-placeholder}
\end{figure}

这也意味着，第五章中的 Orin 独立调试必须把“是否能够以 `cuda` 设备完成推理”作为前置验收条件。若程序在 Orin 上最终退回 CPU，这并不代表平台“没有显卡”，而更可能意味着 JetPack 与 Jetson 版 PyTorch 没有正确对齐。在这种情况下，即使模型可以勉强运行，也不应将其视为满足真机闭环部署要求。

\section{真机部署流程设计}

\subsection{阶段一：Orin 独立调试}

真机部署不应从“让 Lite3 直接走起来”开始，而应先进行 Orin 独立调试。其目标是在不连接 Lite3 的情况下，先验证相机读取、模型加载、CUDA 环境、端到端推理频率和图像保存链路是否正常。这样的顺序设计能够显著降低联调风险，因为一旦高层推理本身尚未稳定，过早接入真机会把问题来源混在一起，难以定位。

在这一阶段，通常需要依次验证以下内容：相机能否稳定读取，NoMaD 权重能否正确加载，统一推理模块能否在 Orin 上输出合理 waypoint，DDIM 参数是否仍处于可接受推理频率，以及 capture 和 session 目录是否能正常生成。对本文来说，这一步的意义不仅是“软件能跑通”，更是确认第四章中形成的统一推理模块真的具备迁移到边缘平台的能力。

在阶段一中，相机标定的判断标准也应一并明确。若相机画面中直线结构基本正常、边缘畸变较小，则可以直接进入独立调试，因为当前代码会自动完成输入尺寸对齐；若相机存在明显桶形畸变或针孔模型偏差，则应先使用离线标定得到相机内参与畸变参数，再在 Orin 端启用桥接层去畸变。为兼顾部署可用性与工程可维护性，本文将相机标定设计为“可选开启”的桥接层功能，而不是强制所有实验场景都先完成标定。

\subsection{阶段一实测结果与分析}

结合 2026 年 4 月 18 日在 Orin 平台上的实际测试输出，本文已完成阶段一独立调试的第一轮实测记录，并将原始命令与终端输出保存到本地结果目录 \texttt{results/deployment/20260418\_orin\_stage1\_validation/orin\_stage1\_validation.md}。这组结果的意义不在于直接给出真机导航成功率，而在于验证高层推理链路、相机链路和部署资产是否已经具备进入阶段二联调的基础条件。

\begin{table}[htbp]
    \centering
    \caption{Orin 阶段一独立调试实测结果}
    \label{tab:orin-stage1-results}
    \begin{tabular}{p{0.20\linewidth}p{0.24\linewidth}p{0.42\linewidth}}
        \toprule
        测试项 & 实测结果 & 结果说明 \\
        \midrule
        相机读取 & $640\times480$，约 $40.0\,\mathrm{ms}$，约 $25.0\,\mathrm{FPS}$ & 说明 USB 相机链路稳定，能够持续提供导航输入图像。 \\
        模型加载 & 2.87\,s，设备为 \texttt{cuda} & 说明 Jetson Orin 集成 GPU 已被正确识别，NoMaD 权重可在板端完成加载。 \\
        基准推理（DDIM-5） & 视觉编码 56.9\,ms；扩散采样 86.7\,ms；端到端 143.6\,ms；约 7.0\,Hz & 说明高层推理模块在 Orin 上已达到可用于后续桥接验证的运行频率。 \\
        完整流水线 & 首轮 2895.5\,ms；后续约 151--157\,ms & 说明系统存在明显冷启动开销，但稳态运行阶段可保持约 6.4--6.6\,Hz。 \\
        Bridge 相机测试 & 连续 30 帧均成功读取 & 说明桥接层相机包装与独立相机测试结果一致。 \\
        Deployment Checklist & 关键文件均存在 & 说明阶段二联调所需的模型、主机、控制器和资产文件链路完整。 \\
        \bottomrule
    \end{tabular}
\end{table}

从结果分析上看，阶段一最重要的结论有三点。第一，相机链路与模型链路已经分别通过验证，因此后续阶段二若出现联调问题，可以优先从网络、UDP 桥接、速度限幅和真机控制链路定位，而不必再把问题归因于“相机无法读取”或“模型根本无法加载”。第二，\texttt{benchmark} 测试给出的约 7.0\,Hz 更适合作为 Orin 高层推理能力的代表值，因为它排除了相机初始化和首轮图优化带来的额外开销；相比之下，\texttt{pipeline} 测试的原始平均值 427.3\,ms 虽然真实，但被第一轮 2895.5\,ms 的冷启动延迟显著拉高，不能直接代表系统进入稳态后的闭环能力。第三，完整流水线在首轮之后基本稳定在 151--157\,ms 区间，说明当前 \texttt{DDIM-5} 配置下系统已经具备进入 Lite3 低速联调的实时性基础。

因此，本文将这组阶段一结果判定为“通过”。这里的“通过”并不意味着真机闭环导航已经完成，而是意味着 Orin 端已经满足以下前置条件：一是图像输入稳定，二是模型能以 \texttt{cuda} 设备运行，三是高层输出 waypoint 数值合理且非零，四是部署资产链路完整。基于这些证据，第五章后续的重点就应从“证明 Orin 能否运行 NoMaD”转向“验证 Orin 与 Lite3 联动时的安全性、控制一致性和闭环有效性”。

\subsection{阶段二：Orin 与 Lite3 联动}

在 Orin 独立调试完成后，系统再进入 Orin 与 Lite3 联动阶段。此时需要完成网络连通、UDP 桥接、起立测试、低速速度测试、交互式导航测试和目标导航测试等环节。其组织原则并非一步到位，而是从低风险动作逐步过渡到完整闭环。

具体而言，先确认 Orin 与 Lite3 运动主机之间的网络和心跳正常；再执行 standup 测试，验证平台能否进入自主控制模式；然后执行低速 twist 测试，验证桥接层能否持续下发速度命令；最后才进入导航主机驱动的探索或目标导航流程。这样的流程之所以必要，是因为对于足式机器人而言，任何一个较底层环节不稳定，都会在高层导航阶段被放大。

\subsection{topomap 与目标图像的真实采集}

与仿真环境不同，真机部署时 topomap 与目标图像需要在真实环境中采集。对于本文研究而言，这一步非常关键，因为 NoMaD 的高层决策依赖于“当前图像与目标图像之间的视觉关系”，如果真实环境中采集的节点图像覆盖不足、视角偏差过大或编号组织混乱，即使高层模型和桥接层本身正确，闭环结果也可能显著退化。

因此，真实 topomap 的采集原则应包括：沿目标路线定间隔记录代表性图像、保持节点顺序与空间顺序一致、尽量避免相邻节点之间过大视角跳变，并为每个导航任务明确当前目标图像或目标节点集合。也正因如此，真机部署不是“把仿真地图直接替换成真实地图”的简单过程，而是需要形成一套可重现的真实环境采集规范。

\begin{figure}[htbp]
    \centering
    \fbox{\rule{0pt}{55mm}\rule{0.86\linewidth}{0pt}}
    \caption{图位占位：真实环境 topomap 采集与目标图像组织示意图（待补）}
    \label{fig:real-topomap-placeholder}
\end{figure}

\section{导航模式、探索模式与图像记录机制}

\subsection{导航模式与探索模式}

基于第四章的统一系统设计，真机端同样支持导航模式与探索模式两类高层运行方式。导航模式下，系统明确接收目标图像或目标节点，并由 NoMaD 生成朝向该目标的局部轨迹；探索模式下，系统弱化明确目标条件，使高层策略依赖当前视觉上下文和训练中学到的先验运动趋势继续前进。由于 goal mask 已经在训练阶段统一了这两种能力\cite{nomad2023}，因此真机端不需要切换完全不同的模型，只需要切换任务模式与目标输入即可。

对真机系统而言，这种统一性十分重要。因为状态机和导航主机不再需要为导航与探索分别维护两套高层逻辑，而只需组织不同任务入口即可。这也使得后续多目标导航能够按照“任务队列负责切换当前目标，NoMaD 负责当前子目标局部轨迹”的思路继续扩展。

\subsection{可视化模式与无可视化模式}

真机部署中另一个需要明确区分的问题是可视化模式与无可视化模式。可视化模式主要用于调试、录屏与实验观察，此时系统会显示实时相机画面、当前目标图像以及必要的状态信息；无可视化模式则面向正式运行，优先保证闭环稳定与计算资源利用，不打开图形窗口，只保留日志与必要图像保存。

这一区分并非使用体验上的细节，而是系统实时性设计的一部分。因为在 Orin 这类边缘计算平台上，GUI 显示和图像刷新会消耗额外资源，从而影响推理周期稳定性。将“是否显示图像”设计为可切换项，意味着系统可以在调试阶段最大化可观测性，而在正式运行阶段最大化稳定性。这也是本文在系统层特别强调“推理模块”和“部署模块”分离的原因之一。

\subsection{实时图像、目标图像与 capture 机制}

当前系统已经支持统一的图像记录机制。真机运行过程中，系统可按 session 目录保存实时前视图像、当前任务目标图像以及交互式 capture 图像。其中，实时图像用于回溯机器人当时真正看到的场景；目标图像用于回溯该阶段任务条件；capture 图像则更多服务于调试与交互式目标选择。通过这三类图像的组合，实验者能够较完整地复原一次真机运行的感知上下文。

这一设计具有明显研究价值。因为对视觉导航系统而言，闭环失败很多时候不是单纯的控制异常，而可能与某一时刻的视觉观测、目标图像偏差或切换节点有关。如果运行过程中没有形成完整图像证据链，后续就很难对失败案例进行可靠分析。换言之，图像记录并非附加功能，而是让第五章实验具有可追溯性的基础。

\begin{figure}[htbp]
    \centering
    \fbox{\rule{0pt}{55mm}\rule{0.86\linewidth}{0pt}}
    \caption{图位占位：真机运行中的实时图像、目标图像与 capture 示例图（待补）}
    \label{fig:real-visualization-placeholder}
\end{figure}

\section{安全机制与实验规范}

\subsection{安全机制设计}

真机实验的首要原则并不是“尽快得到漂亮结果”，而是“确保平台和实验者安全”。因此，本文在部署链路中明确设置了多级安全机制，包括导航主机软件急停、程序中断清理、Lite3 运动主机级急停以及必要时的硬件断电。与此同时，桥接层还负责持续维护底层心跳与速度超时保护，从而保证即使上层程序异常退出，机器人也能在短时间内自动停止运动。

从系统角度看，这些安全机制并不是与 NoMaD 无关的额外功能，而是高层生成式策略能够进入真机的先决条件\cite{diffusionpolicy2023,hwangbo2019anymal,lee2020learning}。因为生成式模型输出本身具有分布性，如果系统没有明确的速度限幅、恢复和急停通道，即使轨迹在大多数时刻合理，也可能在局部异常时对平台造成风险。也正因此，第五章中的平台结论必须同时包含“能否运行”和“能否安全运行”两个维度。

\subsection{实验组织规范}

结合当前部署流程，本文将真机实验组织规范概括为以下几点：初次联调必须使用保守速度限幅；探索与导航任务都应设置合理的最大步数；真机运行必须有人监护；每轮实验前后都应保存 session 日志和关键图像；若系统处于调试阶段，应优先选择可视化模式并保留 capture；若进入正式运行阶段，应优先保证无 GUI 的稳定闭环。

这种规范化组织的价值在于，它使第五章实验从一组临时命令上升为可复现实验流程。对于毕业论文而言，这一点同样重要，因为真机章节如果只有“某次成功运行”而缺少组织原则，就很难形成稳定的研究证据链。

\section{当前完成情况、占位与后续补充方向}

\subsection{当前已可写入正文的内容}

截至当前版本，第五章中已经可以形成稳定正文的内容包括：Orin 作为边缘计算平台的系统定位，Orin 与 Lite3 的软件层次关系，真机桥接层的作用，分阶段部署流程，可视化与无可视化模式的区分，实时图像与目标图像保存机制，以及多级安全机制与实验组织规范。这些内容已经足以构成一章完整的“真机部署方法与实验组织”章节。

\subsection{尚未完成内容的占位说明}

与此同时，仍有若干内容尚未达到最终论文定稿所需的证据强度，因此需要明确占位留空：

\begin{enumerate}
    \item 真机运行照片、现场调试照片、Orin 安装位置照片尚未正式插图，本章已预留图位；
    \item Lite3 真机导航与探索的正式统计结果尚未整理为论文表格，当前只写系统方法，不写绝对性能结论；
    \item 多目标真机任务目前尚未形成完整实验记录，后续可在本章新增“多目标真机实验”小节；
    \item 若后续完成 Tron1 真机部署，则可在本章后半部分扩展“双平台真机对比”内容，目前暂留空位。
\end{enumerate}

\begin{figure}[htbp]
    \centering
    \fbox{\rule{0pt}{58mm}\rule{0.86\linewidth}{0pt}}
    \caption{图位占位：真机实验照片、运行截图与失败案例图（待补）}
    \label{fig:real-pending-placeholder}
\end{figure}

\subsection{建议补充的小节占位}

为了方便后续继续扩展，本章建议预留以下小节作为后续补充入口：

\begin{enumerate}
    \item \textbf{Lite3 真机导航实验结果与分析}：待补正式统计表、成功案例与失败案例；
    \item \textbf{Lite3 真机探索实验结果与分析}：待补探索模式运行轨迹和图像证据；
    \item \textbf{多目标真机任务验证}：待补导航主机在真实任务队列下的执行情况；
    \item \textbf{Tron1 真机部署扩展}：待补双轮足平台的接口适配与实验记录。
\end{enumerate}

\section{本章小结}

本章围绕 Orin 与 Lite3 的真机部署问题，对统一闭环系统如何真正进入真实平台进行了方法层整理。首先，明确了 Orin 在系统中的角色，并说明真机部署为什么需要边缘计算平台、桥接层和分阶段联调流程。其次，围绕导航模式、探索模式、可视化与无可视化模式、实时图像与目标图像保存机制，对真机系统的运行方式进行了细化。再次，从多级急停、速度限幅、心跳保护和实验组织规范等方面，说明了高层生成式导航策略进入真机系统所必须具备的安全条件。

当前阶段，第五章已经能够比较完整地说明本文的真机部署方法、系统接口和实验组织逻辑，但关于正式真机结果、现场照片和多目标任务统计，仍保留若干明确的占位，等待后续实验完成后继续补充。即便如此，本章仍然已经给出了一个清晰结论：第四章建立的统一推理模块、状态机与导航主机架构，已经具备切换到 Orin 与 Lite3 真机平台的结构基础，而本文剩余的工作重点，更多是补充更充分的平台证据，而不是重新推翻整套系统设计。
